% Actualizacion 2024: Logotipos de la UPV


\documentclass[12pt]{article}
\usepackage[left=2.5cm,top=2.0cm,right=2.5cm,bottom=3.0cm]{geometry}
\usepackage[utf8]{inputenc}
\usepackage[spanish]{babel}
%\usepackage[english, american]{babel}
%\usepackage[spanish,es-tabla]{babel}
\usepackage[linguistics]{forest}
\usepackage{amssymb, amsmath, amsbsy} % simbolitos
\usepackage{longtable} % para tablas largas
\usepackage{graphicx}
\usepackage{fancyhdr}
\usepackage{xcolor}
\usepackage{multirow}
\usepackage{listings}
\usepackage{caption}
\usepackage{subcaption}
%\usepackage{parskip}
\usepackage[skip=12pt plus1pt]{parskip}
\usepackage{pdfpages} % Incluir PDF en documento en LATEX
\usepackage{verbatim} % comentarios
\usepackage{algpseudocode}
\usepackage{algorithm}
\usepackage{pdflscape}
\usepackage{multirow}
\usepackage{afterpage}
\usepackage{array,booktabs,ragged2e}
%\newcolumntype{R}[1]{>{\RaggedLeft\arraybackslash}p{#1}}
\newcolumntype{L}[1]{>{\raggedright\let\newline\\\arraybackslash\hspace{0pt}}m{#1}}
\newcolumntype{C}[1]{>{\centering\let\newline\\\arraybackslash\hspace{0pt}}m{#1}}
\newcolumntype{R}[1]{>{\raggedleft\let\newline\\\arraybackslash\hspace{0pt}}m{#1}}

\floatname{algorithm}{Algoritmo}
\renewcommand{\listalgorithmname}{Lista de algoritmos}
\renewcommand{\algorithmicrequire}{\textbf{Entrada:}}
\renewcommand{\algorithmicensure}{\textbf{Salida:}}

% Comentario con respecto a las referencias:
% Por defecto este documento utiliza el formato IEEEtr para formatear las referencias. Si algun asesor requiere formato APA, solo comente la siguiente linea y descomente la linea debajo. Tambien para que las citas funcionen de manera adecuada, el paquete label debe tener las opciones mostradas, o en su defecto, el año en el listado no lo muestra correctamente \usepackage[english, american]{babel}

% Para utilizar el formato de citas IEEE y comentar los dos parrafos siguientes
\usepackage[backend=bibtex,sorting=none]{biblatex}
% Para utilizar el formato APA, sugiero comentar la linea anterior y descomentar las dos proximas lineas
%\usepackage[backend=biber,style=apa]{biblatex}
%\DeclareLanguageMapping{english}{american-apa}

\makeatletter
\DefineBibliographyExtras{spanish}{%
  \setcounter{smartand}{1}%
  \let\lbx@finalnamedelim=\lbx@es@smartand
  \let\lbx@finallistdelim=\lbx@es@smartand
}
\renewbibmacro*{name:delim:apa:family-given}[1]{%
  \ifnumgreater{\value{listcount}}{\value{liststart}}
    {\ifboolexpr{
       test {\ifnumless{\value{listcount}}{\value{liststop}}}
       or
       test \ifmorenames
     }
       {\printdelim{multinamedelim}}
       {\lbx@finalnamedelim{#1}}}
    {}}
\makeatother


% Estas lineas permiten romper los hipervinculos muy largos en las referencias!!!!
\setcounter{biburllcpenalty}{7000}
\setcounter{biburlucpenalty}{8000}
\addbibresource{x.bib} % ARCHIVO DE BIBLIOGRAFÍA


%\usepackage{url}
\usepackage[bookmarks=true,breaklinks=true,bookmarksopen=false,colorlinks=true,linkcolor=blue]{hyperref}
\usepackage[hyphenbreaks]{breakurl}
% Regla que define explicitamente que caracteres rompen los hipervinculos para separar las lineas
%https://es.overleaf.com/11089898rhgykrqyqytx
% Actualiza en automático la fecha de las citas de internet a la fecha de la compilación del documento
\usepackage{datetime}
\newdateformat{specialdate}{\twodigit{\THEDAY}-\twodigit{\THEMONTH}-\THEYEAR}
%\newdateformat{specialdate}{\twodigit{\THEDAY}-\THEYEAR}
\date{\specialdate\today}

\newcommand{\HRule}{\rule{\linewidth}{0.25mm}}


% CONSTANTES NECESARIAS PARA EL DOCUMENTO ---> MODIFIQUEN A SU CRITERIO
\newcommand{\ncarrera}            {Ingeniería en Tecnologías de la Información}
\newcommand{\nasesorinstitucional}{Dr. Said Polanco Martagon}
\newcommand{\NombreAlumno}{
    Mario Coyoy López \\
    Melchor Hernández Díaz \\
    Aldo de Jesús Peña Cuellar \\
    Omar Alejandro Pérez Reyes \\
    Adrian Alejandro Ruíz Márquez \\
    Juan Daniel Torres Colorado
}

%Hombres cambien LA por EL
\newcommand{\elolaNombreAlumno}{el}  
\newcommand{\OA}           {o}  %Hombres cambien A por O
\newcommand{\Matricula}           {1730505}  %SU MATRICULA
\newcommand{\NombreProyecto}{Coche autónomo a escala}
\newcommand{\fechacarta}{26 de Abril de 2021}
\newcommand{\ncuatrimestre}{Mayo-agosto 2021}
\newcommand{\nevalador}{Dr. Hiram Herrera Rivas}
\newcommand{\FechaExposicion}{11 de Agosto de 2021}
\newcommand{\HoraExposionFormatoVenticuatroHoras}{10:00}
\newcommand{\elolaNombreEmpresa}{la}  %Si la empresa es femenino (por ejemplo universidad, usen la) o masculino (el instituto) pongan el
\newcommand{\organismoreceptor}   {Universidad Politécnica de Victoria}

% NOTA: Dos diagonales juntas (\\) indican un saldo de linea. En este caso particular hay 2 (el titulo se ajusta a tres lineas, porque es muy largo. Hacer las adecuaciones pertinentes
\newcommand{\NombreProyectoheader}     {Coche autónomo a escala.}
\newcommand{\nasesorempresaria}   {MSI. José Fidencio López Luna}
\newcommand{\fechaPortada}               {Octubre de 2024}





\newcommand{\separacionCorta}{0.0cm}
\newcommand{\separacionLarga}{0.5cm}

\usepackage[overload]{textcase}
\newcommand{\iemph}[1]{\MakeTextUppercase{#1}}

\pagestyle{fancy}
\headheight 45pt
\fancyhead{} % Clear all header fields
\fancyhead[L]{\includegraphics[height=1.00cm]{UTYP.png}}%
\fancyhead[C]{\begin{center}\NombreProyectoheader\end{center}}%
%\fancyhead[R]{\includegraphics[height=1.5cm]{LogoUPV_2019.png}}%
\fancyhead[R]{\includegraphics[height=1.25cm]{LogoUPV_2023.png}}%
\fancyfoot[R]{\thepage} % Clear all footer fields 
\fancyfoot[C]{}
\fancyfoot[L]{}

\DefineBibliographyStrings{english}{%
  references = {Referencias},% replace "references" with "bibliography"  for `book`/`report`
}

\addto\captionsenglish{%
  \renewcommand{\figurename}{Figura}%
  \renewcommand{\tablename}{Tabla}%
} 

\usepackage{wallpaper}
 
 
%\renewcommand{\figurename}{Figura}
%\renewcommand{\tablename}{Tabla}

 
\begin{document}

%-----------------------------------------------------------------------------------------------------------------
% PAGINA 1 - PORTADA
\setcounter{page}{1}
\pagenumbering{roman}
\thispagestyle{empty}

\begin{center}

\begin{tabular}{cp{5cm}c}
\includegraphics[height=2.25cm]{UTYP.png} & 
& \includegraphics[height=2.25cm]{LogoUPV_2023.png}   \\
\end{tabular}

\Large \textbf{UNIVERSIDAD POLITÉCNICA DE VICTORIA}
\vspace{0.5cm}
\hrule
\vspace{0.1cm} 
\hrule
\vspace{0.5cm}


%\HRule \\[\separacionCorta]
\textbf{\iemph{\NombreProyecto}} \\[\separacionLarga]
%\Large \textbf{TESINA}
%\HRule \\[\separacionLarga]
Sistemas embebidos\\
\textbf{\iemph{\ncarrera}} \\[\separacionLarga]

PRESENTA: \\[\separacionCorta]
%\textbf{\Capitalize{\NombreAlumno}\\[\separacionLarga]
\textbf{\iemph{\NombreAlumno}}\\[\separacionLarga]
%EN CUMPLIMIENTO DE \\[\separacionCorta]
%LA ESTADÍA DE LA CARRERA DE \\[\separacionCorta]


ASESOR \\[\separacionCorta]
\textbf{\iemph{\nasesorinstitucional}} \\[\separacionCorta]

%CO-DIRECTOR \\[\separacionCorta]
%\textbf{\iemph{\nasesorempresaria}} \\[\separacionCorta]

%ORGANISMO RECEPTOR \\[\separacionCorta]
%\textbf{\iemph{\organismoreceptor}} \\[\separacionLarga]

\end{center}
\vspace{4cm}

\begin{flushright}
\iemph{Ciudad Victoria, Tamaulipas, \fechaPortada}
\end{flushright}

\HRule 

\begin{comment}


% En las siguientes 3 paginas, debe incluir una digitalización de calidad de sus respectivas cartaas, no una fotografia toda cucha tomada con su celular de Coppel
% Pagina 1: Digitalización de la carta de presentación (sustituir por una original en el empastado)
\clearpage
\thispagestyle{fancy}
\pagenumbering{roman}
\setcounter{page}{1}
\includepdf[pages={1}]{CartaPresentacion.pdf}

% Pagina 2: Digitalización de la carta de aceptación (sustituir por una original en el empastado)
\clearpage
\thispagestyle{fancy}
\includepdf[pages={1}]{CartaAceptacion.pdf}

% Pagina 3: Digitalización de la carta de liberación (sustituir por una original en el empastado)
\clearpage
\thispagestyle{fancy}
\includepdf[pages={1}]{CartaLiberacion.pdf}

% Pagina 4: Carta de aceptación del ASESOR INSTITUCIONAL(sustituir por una original en el empastado)
\clearpage
\ULCornerWallPaper{1}{Membrete_2023_Sin_Colores_Diabolicos.pdf}
\thispagestyle{empty}

\vspace*{1.5cm}
\large

\begin{center}
\textbf{CARTA DE ACEPTACIÓN DEL DOCUMENTO PARA SU IMPRESIÓN}\\[\separacionLarga]
\end{center}

\begin{flushright}
Cd. Victoria, Tamaulipas a \fechacarta \\[\separacionLarga]
\end{flushright}

\parindent=0mm

\NombreAlumno \\
PRESENTE \\[\separacionCorta]

Le comunico que  el Programa Académico de \ncarrera\ \ le ha otorgado la autorización para la impresión de su Tesina de Estadía Práctica cuyo título es: \\[\separacionCorta]

\begin{center}
\textbf{\NombreProyecto} \\[\separacionLarga]
\end{center}

\begin{center}	
\begin{tabular}{ccc}
\centering
& ATENTAMENTE & \\ 
& & \\
& & \\
& & \\ \hline
& \nasesorinstitucional & \\
& ASESOR INSTITUCIONAL & \\
\end{tabular}
\end{center} 
\vspace{2cm}
c.c.p Director de programa académico


%-----------------------------------------------------------------------------------------------------------------
\clearpage
\ClearWallPaper
\thispagestyle{empty}
\newgeometry{left=1.5cm,top=1.0cm,right=1.5cm,bottom=2.0cm}             
\begin{landscape}

%\afterpage{\restoregeometry}


\begin{tabular}{p{3cm}p{16cm}p{3cm}}
\multirow{4}{*}{\includegraphics[width=3.0cm]{UTYP.png}} &  & \multirow{4}{*}{\includegraphics[width=3.0cm]{LogoUPV_2023.png}} \\
   & \multicolumn{1}{c}{\textbf{EVALUACIÓN DE ESTADÍA}} &  \\ %\cline{2-2}
& \multicolumn{1}{c}{\textbf{Rúbrica para evaluación de la presentación y el reporte de estadía}} & \\ 
& & \\ 
\end{tabular}

\normalsize
\begin{tabular}{p{13cm}R{10cm}}
\multicolumn{1}{l}{Nombre del alumno: \underline{\textbf{\iemph{\NombreAlumno}}}}  & 
Calificación final: \underline{\hspace{3cm}}    \\ %\cline{2-2}
\multicolumn{2}{c}{Periodo: \underline{\textbf{\iemph{\ncuatrimestre}}}}    \\ %\cline{2-2}
& \\
\end{tabular}

\scriptsize 
%\begin{center}
%\resizebox{\linewidth}{!}{%
\begin{tabular}{C{2.0cm}|C{2.0cm}|C{4.5cm}|C{4.5cm}|C{4.5cm}|C{4.5cm}}
\hline
\multirow{2}{*}{Ponderación} & Aspecto a  & Competente & Independiente & Básico Avanzado & No Competente \\
 & Evaluar & 10 & 9 & 8 & 5 \\ 
\hline
40 & Resultados y Actividades & Estrechamente relacionados al perfil de
egreso de su programa académico & Parcialmente relacionados al perfil de
egreso de su programa académico & Escasamente relacionados al perfil de egreso de su programa académico & Escasamente relacionados al perfil de egreso de su programa académico \\  \hline

30 & Exposición de las actividades de la estadía  & Detalladas y sustentadas con respecto a los resultados que se obtuvieron & Detalladas y sustentadas parcialmente con respecto a los resultados que se obtuvieron & Detalladas parcialmente con respecto a los resultados que se obtuvieron & Detalladas escazamente con respecto a los resultados que se obtuvieron \\  \hline 

10 & Material visual Lenguaje
verbal & Uso el lenguaje y la terminología
apropiadas; El material visual está organizado,
adecuado y suficiente
 & 
Uso el lenguaje y la terminología apropiadas
El material visual está parcialmente
organizado y es suficiente
& 
Uso el lenguaje y la terminología son parcialmente apropiadas; El material visual está parcialmente
organizado y es suficiente
 & 
Uso el lenguaje y terminología es inapropiado; El material visual no está organizado y es insuficiente
 \\  \hline
10 & Exposición en Idioma Inglés &

Pronunciation is clear so language is easily understood (2.5) Uses fluent connected speech, occasionally disrupted by search for correct form of expression (2.5) Uses topic related vocabulary without problems (2.5) Responds to questions using varied and descriptive vocabulary and language structures (2.5) & 

Pronunciation is understandable, but there are slight errors (2.25)
Speech is connected but frequently disrupted by search for correct form of
expression (2.25) Uses some topic related vocabulary sufficient to communicate ideas (2.25) Responds to questions using simple but accurate vocabulary and language structures (2.25 & 

Pronunciation is understandable most of the time, marked native accent and many errors
(2) Speaks with simple sentences, sometimes not connected, but is understood (2) Uses basic vocabulary to communicate ideas
(2) Partly responds to simple questions, with limited vocabulary and language structures (2)  &

Pronunciation makes language very difficult to understand (1) Uses one-word/two-word utterances (1) Unable to communicate ideas due to lack of vocabulary (1) Uses isolated words or sentence fragments to respond to questions (1) 
\\  \hline

5 & Respuesta a los cuestionamientos de los evaluadores & 
Clara y satisfactoria & 
Clara y parcialmente satisfactoria & 
Clara e insuficiente & 
Confusa e insuficiente 
\\  \hline
5 & Autorización de tesina en tiempo y forma & 
Presenta en tiempo y forma 
\tikz[overlay, remember picture,anchor=base]  \node (Mark){}; 
 & 
Presenta en tiempo y forma con la mayoría de requerimientos solicitados & 
Presenta en tiempo y con algunas limitantes  de los requerimientos solicitados. &
Presenta fuera de tiempo y con los mínimos requerimientos solicitados.
 \\  % Una linea en Blanco para poner la marca del ASESOR INSTITUCIONAL
&
&
%\tikz[overlay, remember picture,anchor=base] \node (Center){};

&
 \\  \hline

\end{tabular}

%\begin{tikzpicture}[remember picture, overlay, note/.style={rectangle callout, fill=#1}]
%\node [note=red!80, callout absolute pointer={(Mark)}] at (Center) {COMPETENTE!};
%\end{tikzpicture}

%}
%\end{center}

\normalsize


%\begin{center}
%\begin{tabular}{ccc}
%\includegraphics[scale=0.10]{FirmaMANM2.png} & & \\
%\hline 
%\nasesorinstitucional & \nevalador & \\
%ASESOR INSTITUCIONAL & EVALUADOR & EVALUADOR INGLÉS \\
%\end{tabular}
%\end{center}

\begin{center}
\begin{tabular}{cp{1cm}cp{1cm}c}
%\includegraphics[scale=0.10]{ArchivoFirmaASESOR.png} & & & & \\
 & & & & \\
  & & & & \\
\hline 
\nasesorinstitucional & & \nevalador & & \\
ASESOR INSTITUCIONAL & &EVALUADOR & &EVALUADOR DE INGLÉS \\
\end{tabular}
\end{center}


\end{landscape} 
\restoregeometry

%-----------------------------------------------------------------------------------------------------------------
% Pagina 5: Registro de Evaluación de Estadía
\clearpage
\pagestyle{empty}
\ULCornerWallPaper{1}{Membrete_2023_Sin_Colores_Diabolicos.pdf}

\vspace*{0.5cm}



\large
\begin{center}
\textbf{REGISTRO DE EVALUACIÓN DE EXPOSICIÓN DE ESTADÍA}
\\[\separacionLarga]
\end{center}


Siendo las \HoraExposionFormatoVenticuatroHoras \ \ horas del día \FechaExposicion, \elolaNombreAlumno\ \ alumn\OA\ \ \textbf{\NombreAlumno}, del programa académico \textbf{\ncarrera}, con matricula \textbf{\Matricula}, presentó la exposición de la estadía realizada durante el cuatrimestre \textbf{\ncuatrimestre}, en \elolaNombreEmpresa\ \ \textbf{\organismoreceptor}, con el proyecto titulado \textbf{\NombreProyecto}.\\

Una vez concluido el proceso de evaluación, y con base a la rúbrica establecida para éste propósito, se determina que la calificación de la estadía es \underline{\hspace{2cm}}. %\hrulefill.

\begin{center}
\begin{tabular}{ccc}
& & \\
& & \\
& & \\
%& \includegraphics[scale=0.10]{FirmaAsesorPro.png} & \\
\hline 
%& \underline{\hspace{8cm}}& \\
& \nasesorinstitucional & \\
& ASESOR INSTITUCIONAL & \\
& & \\
& & \\
& & \\
%& \underline{\hspace{8cm}}& \\
\hline 
& \nevalador & \\
& EVALUADOR & \\
& & \\
& & \\
%& & \\
\hline 
& & \\
& EVALUADOR DE INGLÉS & \\
\end{tabular}
\end{center}

\normalsize

%-----------------------------------------------------------------------------------------------------------------
% PAGINA 4 - AGRADECIMIENTOS
\clearpage
\ClearWallPaper
\pagestyle{fancy}
\section*{\centering Agradecimientos}
\addcontentsline{toc}{section}{Agradecimientos}
%\input{Agradecimientos.tex}

%-----------------------------------------------------------------------------------------------------------------
% PAGINA 5 - RESUMEN EN ESPAÑOL

\clearpage
\section*{\centering Resumen}
\addcontentsline{toc}{section}{Resumen}
%\input{Resumen.tex}
\textbf{Palabras clave:} Tecnología, Sistema web, Desarrollar, Implementar, Empresas, Modulo.

%-----------------------------------------------------------------------------------------------------------------
% PAGINA 6 - RESUMEN EN INGLES

\clearpage
\section*{\centering Summary}
\addcontentsline{toc}{section}{Summary}
%\input{Summary.tex}
\textbf{Keywords}: Technology, Web system, Develop, Implement, Companies, Module.

%-----------------------------------------------------------------------------------------------------------------
\end{comment}

% PAGINA 7 - INDICE

\clearpage
\addcontentsline{toc}{section}{Índice}
\renewcommand\contentsname{Índice}
\tableofcontents

%-----------------------------------------------------------------------------------------------------------------
% CAPITULOS

\clearpage
\pagenumbering{arabic}
\setcounter{page}{1}
%\input{Capitulo1.tex}

\clearpage
%\input{Capitulo2.tex}

\clearpage
\section{Introducción}
La conducción  autónoma ha sido uno de los mayores retos y con creciente interés en los últimos años, sobretodo cuando esta implica transformar la forma en que nos desplazamos. En el año 1939, se realizó el primer coche eléctrico capaz de ser controlado por circuitos integrados, un avance para que años más tarde, en 1994 ocurriera un hito en la historia automotriz. Dos coches condujeron de manera autónoma por más de dos mil kilómetros en una autopista de París, donde la intervención humana no estaba exenta, pero era muy reducida.

Hoy en día, existen prototipos, investigaciones y Sistemas Avanzados de Asistencia a la Conducción o \textit{Advanced Driver Assistance System} (ADAS), que incluso ya están en el mercado, pero aún están lejos de liderarlo, pues para ello primero se debe lograr una normalización sobre los vehículos.

Esto se ha convertido en una carrera tecnológica para los fabricantes y las grandes marcas, aún así, el problema que buscan resolver a través de la conducción autónoma, sigue teniendo el mismo enfoque, reducir los accidentes causados por los errores humanos, mejorar el trafico, facilitar la movilidad de aquellas personas con discapacidad o limitaciones físicas y que con ello sea un transporte más ecológico, seguro y eficiente. 

En las grandes urbes es cada vez más común ver modelos de vehículos que, en mayor o menor medida, empiezan a ser autónomos: realizan la maniobra de estacionamiento por si solos, circulan sobre avenidas altamente transitadas o en otros casos reaccionan ante una posible situación de choque o atropello, por otro lado, aún falta tiempo para que se comercialice un vehículo totalmente autónomo, que pueda navegar sin intervención humana y ante cualquier situación.

No obstante, pesar de las grandes compañías e investigaciones que abordan este campo, aún existen diversas limitaciones tecnológicas y legales que impiden la adopción de los coches autónomos en su totalidad. Uno de los principales problemas es la infraestructura urbana. Las ciudades, están aún lejos de tener el entorno adecuado para soportar su normalización, pues su aplicación no es trivial.

El ambiente por donde el vehículo circula, es muy diverso, incluye variaciones climatológicas (días soleados, nublados, lluvia, nieve, etc.), cambios en el terreno (pavimento, asfalto, tierra, piedras, etc.), obstáculos en el camino, señalamientos de tránsito, rutas aún no existentes en mapas digitales; lo que vuelve este problema sumamente complejo para la percepción de su entorno.

Desde el punto de vista tecnológico, la inteligencia artificial (IA) ha transformado múltiples sectores, y se ha exponenciado como una solución ante los problemas actuales, sin embargo, a pesar de sus avances, la IA enfrenta limitaciones cuando se trata de reaccionar antes situaciones impredecibles y de procesar información desconocida. Este desafío se vuelve particularmente crítico en contextos donde el tiempo de reacción se traduce en decisiones con altas consecuencias. Si bien existen sensores que nos permiten obtener información en tiempo real, la latencia, se vuelve un obstáculo para su efectividad.

Además, el tiempo que tarda un dato en viajar hacia su destino para ser procesado, y la posterior respuesta que permite al sistema actuar, son factores críticos en los que cada segundo cuenta. Esto requiere que tanto la recopilación, el procesamiento y la toma de decisiones se encuentre optimizada al origen de los datos, es decir, la computación de borde o \textit{edge computing}.

El \textit{edge computing}, es un enfoque clave para resolver este problema, ya que permite procesar los datos lo más cerca posible de donde se generan, en este caso, dentro del propio vehículo. Al procesar la información localmente, reduce la latencia y mejora la capacidad de respuesta del sistema. 

Es por ello que los sistemas embebidos juegan un papel crucial, ya que estos dispositivos compactos y de bajo consumo pueden realizar tareas de procesamiento local, como el reconocimiento de imágenes, necesario para que un vehículo autónomo navegue.

Sin embargo, surge una pregunta elemental: ¿pueden los sistemas embebidos cumplir con estas demandas de procesamiento de datos y toma de decisiones? Estos dispositivos, aunque eficientes en consumo energético y de tamaño compacto, presentan limitaciones de capacidad de cómputo en comparación con soluciones centralizadas o en la nube.

En este contexto, este proyecto plantea el diseño y desarrollo de un coche autónomo a escala, utilizando una placa ESP32 como plataforma de procesamiento. La hipótesis central de este trabajo es que, a pesar de las limitaciones inherentes a los sistemas embebidos, la ESP32 posee la capacidad suficiente para ejecutar tareas de reconocimiento de imágenes y control en tiempo real, necesarias para la conducción.

\clearpage
\section{Antecedentes Teóricos}
La evolución de la conducción autónoma, se trata de uno de los mayores avances tecnológicos de este siglo, en el que el vehículo se desplaza por sí solo sin necesidad de esfuerzo humano para conducirlo (bajo ciertas limitaciones) o incluso mantenerlo bajo control. Esto también implica sistemas de transporte autónomos y robots. Por tanto, implica que estos vehículos estén dotados de sensores, actuadores y sistemas de control a través de los cuales se vuelven capaces de moverse de forma autónoma.

Los cinco niveles de conducción autónoma, que van desde el control manual total (Nivel 0) hasta la autonomía completa (Nivel 5), marcan la evolución de esta tecnología. Actualmente, los fabricantes están avanzando en la comercialización de vehículos de Nivel 3, como los de Mercedes, que operan de manera autónoma bajo ciertas condiciones en autopistas, mientras que Tesla sigue perfeccionando su sistema Full Self-Driving con un enfoque más agresivo hacia el Nivel 4 \cite{mlearning2021}.

Las tecnologías comunes en los sistemas de conducción autónoma incluyen equipos de control al que se conectan diversos sensores, radares, cámaras, etc., que observan regularmente el entorno desde adentro hacia afuera en busca de todos los objetos, vehículos, personas, entre otros elementos en su campo de visión. Con la aparición de tecnologías como TinyML, que implementa modelos de aprendizaje automático en dispositivos de baja potencia, la oportunidad de diseñar prototipos mejorados y más accesibles ha proporcionado una nueva serie de oportunidades para la investigación y el desarrollo.

El aumento de la demanda de integrar dispositivos de señales de entrada de bajo consumo y la creación de modelos de ML eficientes en energía y memoria han facilitado el desarrollo de TinyML. TinyML tiene como objetivo la implementación del aprendizaje automático y las prácticas en dispositivos de señal, así como en la vanguardia de los marcos de IoT. Esto permite el procesamiento de señales y la inteligencia en dispositivos IoT, evitando la necesidad de cargar datos en servicios externos para su procesamiento. Al implementar de manera efectiva TinyML en dispositivos de Internet de las cosas, se pueden lograr niveles más altos de privacidad y eficacia, al mismo tiempo que se reducen los costos. Además, en situaciones donde la conectividad no es necesaria, el principal atractivo de TinyML radica en su enfoque en el análisis.


Los últimos avances en vehículos autónomos se han centrado en la optimización de la toma de decisiones utilizando enfoques como el aprendizaje reforzado basado en procesos de decisión de Markov. Este marco matemático permite que los vehículos elijan acciones óptimas y ajusten sus trayectorias en tiempo real a través de recompensas por sus decisiones. Fabricantes como Tesla, Waymo y Nvidia están implementando estrategias que combinan sensores como LiDAR, cámaras y radar para mejorar la percepción del entorno. Además, el desarrollo de algoritmos de aprendizaje profundo extremo a extremo ha permitido la interpretación directa de imágenes de cámaras en comandos de conducción, aunque la complejidad del entorno sigue siendo un reto para la predicción de situaciones imprevistas. Este avance se considera un paso hacia la inteligencia artificial general, donde los vehículos podrían deducir nuevos contextos de conducción, una capacidad que es aún limitada pero clave para la evolución futura de estos sistemas hacia una verdadera inteligencia artificial general \cite{mlearninglabAvances}.

\textbf{Detección de Objetos en Vehículos Autónomos} 
\\
TinyML puede tener aplicaciones generalizadas en la industria de vehículos autónomos, ya que estos vehículos se pueden utilizar de diferentes maneras, incluido el seguimiento humano, fines militares y aplicaciones industriales. Estos vehículos tienen el requisito principal de poder identificar objetos de manera eficiente cuando se busca el objeto. 
A partir de ahora, los vehículos autónomos y la conducción autónoma son una tarea bastante compleja, especialmente cuando se desarrollan vehículos mini o de tamaño pequeño. Los desarrollos recientes han demostrado potencial para mejorar la aplicación de la conducción autónoma para mini vehículos mediante el uso de una arquitectura CNN y la implementación del modelo sobre el MCI GAP8. 
En los últimos cinco años se ha estado trabajando más para la evolución y mejora de los vehículos con conducción autónoma, entre ellos el uso de la inteligencia artificial para la toma de decisiones en tiempo real, gran parte de estos avances se centran en la detección y reconocimiento de obstaculos, peatones, vehículos, entre otros, para garantizar la funcionalidad de los vahiculos autónomo.
Uno de los campos claves en esta evolución ha sido el desarrollo de algoritmos avanzados de visión por computadora, y redes neuronales como YOLO.  
A lo largo del tiempo, se han desarrollado una gran variedad de automóviles autónomos utilizando Arduino, los cuales de manera general abarcan desde el diseño físico, la implementación electrónica y la programación de software. Algunos trabajos preexistentes han explorado la detección de carriles y señales de tráfico mediante algoritmos de procesamiento de imágenes en tiempo real, utilizando cámaras y sensores acoplados al vehículo. Estos enfoques han permitido que los vehículos autónomos puedan identificar la delimitación de las vías de manera precisa, incluso en condiciones de baja visibilidad o en carreteras mal señalizadas.


\subsection{Marco Teórico}

\textbf{Vision Por Computadora}

La visión por computadora es un campo de la inteligencia artificial que permite a los vehículos autónomos 'ver' el entorno. Los algoritmos de detección de objetos, como YOLO (You Only Look Once) y Faster R-CNN, han sido esenciales para que los vehículos puedan identificar y clasificar objetos, como otros autos, peatones y señales de tráfico, de manera rápida y precisa. Esta metodología utiliza redes neuronales convolucionales para analizar imágenes capturadas por cámaras en tiempo real \cite{Vision}.

\textbf{Teoría de la conducción autónoma}

Es la base para regular el comportamiento de sistemas autónomos como los vehículos con conducción autónoma. Esta teoría incluye el uso de sistemas de control en bucle cerrado (feedback control), donde el vehículo responde a su entorno en tiempo real mediante sensores y algoritmos. Los sistemas de control PID (Proporcional-Integral-Derivativo) son ampliamente utilizados para mantener estabilidad y precisión en la conducción autónoma \cite{teoria}.

\textbf{Machine Learning}

El machine learning ha permitido grandes avances en la capacidad de los vehículos autónomos para aprender de datos y mejorar su rendimiento en la conducción. A través de redes neuronales profundas, los vehículos pueden analizar vastas cantidades de datos sensoriales y tomar decisiones inteligentes \cite{ML}.

\textbf{Redes Neuronales Artificiales}

Las Redes Neuronales Artificiales han sido aplicadas para resolver problemas de clasificación y predicción en vehículos autónomos. Estas redes imitan el funcionamiento del cerebro humano y permiten que el vehículo aprenda de experiencias anteriores \cite{RN}.


\textbf{Algoritmos de navegación} 

Para el control de la navegación de un robot móvil existen infinidad de alternativas. El SLAM (Simultaneous Localization and Mapping) es un algoritmo clave para la navegación autónoma, ya que permite que los vehículos construyan un mapa de su entorno mientras se localizan dentro de él. Este algoritmo se basa en una combinación de sensores y modelos matemáticos para proporcionar información precisa sobre la posición del vehículo, lo que es esencial para la planificación de rutas y la toma de decisiones en tiempo real \cite{An}.

\subsection{Estado del arte} 

\textbf{Introducción}

El avance de la tecnología de vehículos autónomos ha sido notable en los últimos años, impulsado por el desarrollo de sistemas de aprendizaje automático y la integración de técnicas de computación de borde, como TinyML. Este estado del arte examina las tendencias, métodos y aplicaciones más relevantes en el ámbito de la conducción autónoma, incluyendo el uso de aprendizaje profundo, prototipos, sensores, visión artificial, desafíos regulatorios, y algoritmos avanzados. A continuación, se presentan investigaciones y enfoques recientes que contribuyen a estos avances, basados en estudios y artículos actuales.

\textbf{Aprendizaje Profundo en la Conducción Autónoma}

En su artículo “El impacto del aprendizaje profundo en el desarrollo de vehículos autónomos”, Gary Silberg destaca cómo el aprendizaje profundo permite que los vehículos tomen decisiones sin que los ingenieros tengan que programar cada posible escenario. La recopilación y análisis de datos son fundamentales para que los vehículos se adapten a situaciones impredecibles. Se discuten dos enfoques principales:
\begin{itemize}
    \item \textbf{Abstracción semántica:} Utilizado en situaciones donde la comprensión del contexto es crucial. El modelo empleado son redes neuronales convolucionales (CNN), logrando una precisión del 92\% en la detección de obstáculos.

    \item \textbf{Aprendizaje de punta a punta:} Este enfoque simplifica el proceso al utilizar grandes volúmenes de datos, aunque enfrenta retos en situaciones no vistas previamente. También se basa en CNN, pero requiere mayor capacidad de datos \cite{silberg2024}.

\end{itemize}

\textbf{ Desarrollo de Prototipos de Vehículos Autónomos}
Un proyecto de fin de máster en la Universidad de Valladolid describe el desarrollo de un vehículo autónomo a escala que emplea algoritmos básicos en la plataforma Arduino. El prototipo incluye funciones como el seguimiento de líneas y la evitación de obstáculos, lo que lo convierte en una herramienta valiosa para la evaluación temprana de tecnologías autónomas. Se utiliza un algoritmo de seguimiento de líneas basado en visión artificial, destacando su simplicidad y utilidad en entornos controlados \cite{valla2023}.

\textbf{TinyML y su Aplicación en Dispositivos Autónomos}

Un artículo de Unite.ai examina la creciente relevancia de TinyML, una tecnología que permite ejecutar modelos de aprendizaje automático en dispositivos de bajo consumo energético. Su aplicación en vehículos autónomos permite realizar procesamiento local sin depender de la nube, mejorando la eficiencia y privacidad. Un ejemplo es el uso de TinyML para el procesamiento de imágenes en drones autónomos, logrando evitar obstáculos en tiempo real. El modelo utilizado es TinyML basado en redes neuronales convolucionales (CNN), optimizado para dispositivos de bajo consumo \cite{unite2023}.

\textbf{Sensores y Tecnología en la Navegación Autónoma}

Los avances en sensores como cámaras, LIDAR y radares son fundamentales para la percepción del entorno en vehículos autónomos. Estos sensores trabajan en conjunto para ofrecer una visión completa del entorno y mejorar la interacción segura con el mismo, incluso en escenarios complicados como tráfico urbano o áreas rurales. El sistema de fusión de sensores utilizado, combinado con algoritmos avanzados de percepción, ha alcanzado una precisión del 95\% en la detección de peatones y vehículos, lo que resalta la capacidad de estos sistemas para mejorar la seguridad en situaciones complejas \cite{programarfacil2023}.

\textbf{Sistemas de Visión Artificial en UAVs}

El desarrollo de un sistema de visión artificial para cuadricópteros, que utiliza técnicas de aprendizaje profundo para detectar y rastrear objetos en movimiento, demuestra la versatilidad de los vehículos autónomos en entornos aéreos. Este proyecto utiliza simulaciones con ROS y técnicas basadas en el modelo SSD (Single Shot Multibox Detector) para mejorar la precisión del rastreo, alcanzando una precisión del 90\% en la detección y seguimiento de objetos. Esto resalta la importancia de las simulaciones en el desarrollo de sistemas autónomos, ya que permiten probar y ajustar algoritmos en situaciones controladas antes de implementarlos en vehículos reales \cite{escolar2022}.

\textbf{Retos y Futuro de los Vehículos Autónomos}

A pesar de los avances, el camino hacia la autonomía total está plagado de desafíos, incluidos los problemas relacionados con la regulación, la seguridad y la aceptación pública. Se anticipa que para 2025, se lanzarán modelos con niveles avanzados de autonomía, pero será esencial preparar a los conductores y a las infraestructuras para este cambio. Además, la adopción de estos vehículos enfrenta desafíos en términos de aceptación pública, aspectos éticos sobre la toma de decisiones automáticas en situaciones críticas, y las regulaciones que varían significativamente entre países \cite{atrain2023}.

\textbf{Algoritmos y Desafíos en la Conducción Autónoma}

El desarrollo de algoritmos como matrices de decisión, reconocimiento de patrones y aprendizaje reforzado ha sido clave para el progreso de los vehículos autónomos. Estos algoritmos permiten que los vehículos identifiquen su ubicación, determinen su destino y tomen decisiones críticas como frenar o girar. Por ejemplo, Tesla utiliza matrices de decisión en su sistema de autopilotaje, mientras que Waymo emplea una combinación de reconocimiento de patrones y aprendizaje reforzado para adaptarse a las cambiantes condiciones del tráfico urbano. Estos sistemas han logrado una precisión del 93\% en condiciones de tráfico urbano, lo que subraya la eficiencia de los algoritmos en la conducción autónoma \cite{digiltea2024}.

\textbf{Avances en Aprendizaje Profundo para la Conducción Autónoma}

En el artículo titulado "Deep Learning for Autonomous Vehicles: State-of-the-Art and Future Trends" (2021), se realiza un análisis exhaustivo sobre las técnicas de aprendizaje profundo aplicadas a vehículos autónomos. Los autores examinan arquitecturas como las Redes Neuronales Convolucionales (CNN), las Redes Neuronales Recurrentes (RNN) y las Redes Neuronales de Memoria a Largo Plazo (LSTM), que desempeñan un papel crucial en los sistemas de percepción y toma de decisiones de los vehículos. Además, destacan cómo estas arquitecturas permiten una detección de objetos con una precisión del 95\% y una tasa de falsos positivos del 5\%. También se menciona un tiempo de respuesta del sistema de 200 ms, lo que refleja la eficiencia de estas técnicas en entornos dinámicos \cite{silberg2021deep}.

\textbf{Avances en Técnicas de Aprendizaje Profundo en la Conducción Autónoma}

En el artículo "A Comprehensive Review on Autonomous Navigation, se ofrece una visión integral de las técnicas de aprendizaje profundo aplicadas a la conducción autónoma, cubriendo aspectos desde la percepción hasta la planificación y el control de los vehículos. Los autores subrayan los avances en el uso de Redes Neuronales Convolucionales (CNN) para el reconocimiento de imágenes y la segmentación, así como el Aprendizaje por Refuerzo Profundo (DRL) para la toma de decisiones en entornos complejos. También se exploran aplicaciones de Redes Generativas Antagónicas (GAN) para generar datos sintéticos que entrenen estos sistemas autónomos. En cuanto a las métricas de rendimiento, se reporta una precisión del 93\% en la segmentación de imágenes, una tasa de éxito del 90\% en la navegación autónoma y un tiempo de procesamiento por imagen de 150 ms \cite{nahavandi2022comprehensive}.

\textbf{Comparativa de Modelos de Aprendizaje Profundo en Vehículos Autónomos}

A continuación, se presenta una tabla comparativa que resume los artículos analizados en el contexto de la tecnología de vehículos autónomos. La tabla incluye el título de cada artículo, el modelo de aprendizaje profundo utilizado y la precisión lograda en las respectivas aplicaciones
\begin{table}[ht]
    \centering
    \small % Cambia a \footnotesize si necesitas más reducción
    \begin{tabular}{|p{5cm}|p{5cm}|p{2cm}|} % Ajusta los anchos según sea necesario
        \hline
        \textbf{Título del Artículo} & \textbf{Modelo Utilizado} & \textbf{Precisión} \\ \hline
        El impacto del aprendizaje profundo en el desarrollo de vehículos autónomos & Redes Neuronales Convolucionales (CNN) & 92\% \\ \hline
        Sensores y Tecnología en la Navegación Autónoma & Fusión de sensores y algoritmos & 95\% \\ \hline
        Sistemas de Visión Artificial en UAVs & SSD (Single Shot Multibox Detector) & 90\% \\ \hline
        Deep Learning for Autonomous Vehicles: State-of-the-Art and Future Trends & CNN, RNN y LSTM & 95\% \\ \hline
        Autonomous Driving: A Survey of Deep Learning Techniques & CNN, GAN y Aprendizaje por Refuerzo Profundo (DRL) & 93\% \\ \hline
    \end{tabular}
    \caption{Comparativa de Modelos de Aprendizaje Profundo en Vehículos Autónomos}
    \label{tab:comparativa_modelos}
\end{table}

La investigación y desarrollo en el campo de los vehículos autónomos y tecnologías como TinyML han demostrado un avance significativo en los últimos años. Las aplicaciones son amplias y variadas, desde vehículos terrestres hasta aeronaves no tripuladas, y el futuro promete un crecimiento aún mayor en este campo. La integración de inteligencia artificial en dispositivos cotidianos se vuelve crucial, lo que plantea la necesidad de una formación continua y un enfoque regulador adaptado a estos avances.




\clearpage
\section{Motivación}
La creciente problemática de los accidentes de tráfico y la congestión vial en las grandes y pequeñas urbes impulsa la búsqueda de soluciones. La conducción autónoma busca mejorar la seguridad vial y optimizar la movilidad urbana. Reduciendo la intervención humana que es causante de diversos errores que terminan en accidentes automovilísticos.
Diversas entidades, investigadores y empresas han abordado el problema de la conducción autónoma desde diferentes perspectivas, es por ello que cada día es más común observar coches que cuentan con algún dispositivo de asistencia que permite escalar en los niveles de autonomía. A diferencia de los grandes recursos materiales y tecnológicos que cuentan estas entidades, es relevante aportar en el desarrollo de algoritmos de percepción, planeación y control para vehículos sin conductor donde el procesamiento ocurra en dispositivos embebidos, concretamente el uso de la placa ESP32, donde  es desafiante abordar un problema complejo donde la capacidad del hardware es limitada.


\section{Justificación}
Contribuir al desarrollo y la investigación en el campo de los coches autónomos con un enfoque específico en la integración de dispositivos embebidos. Al centrarnos sobre el uso de la ESP32, se plantea  una solución de \textit{edge computing} que permite el procesamiento de datos en el propio vehículo, agilizando los tiempos de respuesta y la toma de decisiones en situaciones criticas. Este enfoque permite evaluar el rendimiento y funcionamiento con dicho recurso, sin la necesidad de invertir en hardware costoso. No solo servirá como una investigación educativa, si no que al abordar problemas complejos de manera directa y práctica, se traduce en proponer un enfoque que pueda ser aplicado en entornos industriales, promoviendo un futuro donde los vehículos autónomos no sean un lujo sino una realidad accesible para mejorar la vida urbana.

\clearpage
\input{Capitulo6.tex}

\clearpage
\section{Descripción del problema}

La conducción autónoma, como su nombre lo indica, es toda conducción de vehículos que se realice sin conductor. A día de hoy existen ejemplos de autos autónomos, tal como es el caso de los vehículos Tesla que usan el sistema Autopilot, y que sin embargo, su conducción requiere de asistencia y supervisión de un conductor hasta cierto punto. Esto último denota que  se requiere obligatoriamente de un ser humano, al menos por el momento, para realizar determinadas situaciones que requieran del control parcial o total del usuario.

Sin embargo, hay que tener en cuenta que este concepto de 'Conducción autónoma' es muy poco descriptivo por sí mismo, dado que existen ciertos niveles de conducción autónomas.
Como se intuye por lo anterior, para lograr una conducción autónoma, ya sea en cualquiera de sus niveles propuestos, se requiere de hardware y software sofisticados. En este caso, para el proyecto propuesto, el cual es un coche autónomo con ayuda de Machine Learning (TinyML), el hardware con el que se cuenta es limitado, pues se hará uso de microcontroladores, específicamente de una placa ESP32, la cual en cualquiera de sus modelos, su capacidad de procesamiento es bastante limitada, pero que por sí sola cuenta con características que la vuelven bastante útil, como la integración de Wi-Fi u otras características.

Para lograr este reto, la manera más simple es usar librerías de Machine Learning específicamente diseñadas para este tipo de microcontroladores, como es el caso de Arduino Eloquent. Específicamente, para realizar la conducción autónoma, se harán uso de cámaras y de visión por computadora, en donde el vehículo debe de seguir una ruta marcada y deteniéndose ante obstáculos. Para ello, la idea es tener una cámara que permitirá acceder al espectro visual, a partir de eso, el propio microcontrolador deberá tomar una decisión. Este análisis se realizará por medio de algoritmos de visión por computadora, como la detección de bordes, reconocimiento de patrones y/o segmentación de objetos (detección).

Adicionalmente, se deben considerar condiciones operativas variadas. Por ejemplo, las cámaras deben estar despejadas para garantizar una visibilidad óptima; en caso contrario, el vehículo debe ser capaz de detenerse automáticamente y alertar al usuario sobre la situación

Así, el desafío no solo radica en implementar la tecnología adecuada, sino también en garantizar que el prototipo pueda funcionar de manera confiable y segura en condiciones del mundo real. 





\clearpage
\section{Alcances y Limites}
Para este proyecto, los programas elaborados y el equipo utilizado para la detección de objetos en vehículos autónomos serán explícitamente para la conducción autónoma. Se planea que con el uso de la tecnología TinyML con redes neuronales artificiales se logren obtener resultados precisos, cumpliendo así con el objetivo del proyecto previamente antes mencionado. Además, se busca llevar a cabo el desarrollo de un coche con capacidades autónomas, es decir, que el coche tenga la funcionalidad de detectar objetos estáticos que lleguen a interrumpir la trayectoria en línea recta del carril diseñado. Por último, se pretende realizar distintos análisis de las pruebas realizadas durante el proyecto, en la cuál, se busque una mejoría de la precisión del modelo entrenando, notando que el resultado sea óptimo. No obstante, no se implementará una interfaz en la cual el usuario pueda interactuar con la ESP-WROOM-32, debido a que realizar dicho proceso consumirá más recursos necesarios, produciendo una saturación innecesaria afectando el rendimiento.

Algunas de las limitaciones con las que el desarrollo del proyecto tendrá son principalmente en las capacidades de la ESP-WROOM-32. De las cuales, algunas de las características más importantes y destacables con las que se cuenta, son:
\begin{itemize}
    \item \textbf{Almacenamiento interno}: 4MB.
    \item \textbf{Memoría volátil}: 520KB de SRAM.
    \item \textbf{Alimentación}: 3.3V (aunque la placa de desarrollo puede aceptar 5V a través de USB).
    \item \textbf{Temperatura de Operación}: Menores a 85°C. El rendimiento puede verse afectado cuando alcanza condiciones extremas de calor.
\end{itemize}
Estas limitaciones llegarán a limitar el rendimiento obtenido de parte del uso de la tecnología TinyML. A pesar de que EloquentArduino permite facilitar la implementación de los modelos de aprendizaje automático, estos suelen requerir de una optimización significativa para ajustarse a este espacio, lo que limita el tamaño y la complejidad del modelo que se llegue a diseñar, restándole precisión en la que el modelo ha sido entrenada.

Además, otra de las limitaciones que se tendrá es el tiempo. Es decir, el proyecto únicamente contará con un tiempo de desarrollo aproximado de 9 semanas. Esta limitación supone una restricción debido a que la colecta de imágenes que servirán para el conjunto de datos (dataset) para entrenar al modelo puede llegar a ser largo y requerir de una gran cantidad de recursos. 

\begin{comment}

\clearpage
\addcontentsline{toc}{section}{Índice de figuras}
\renewcommand\listfigurename{Índice de figuras}

\listoffigures

\clearpage
\addcontentsline{toc}{section}{Índice de cuadros}
\renewcommand\listtablename{Índice de cuadros}
\listoftables

\clearpage
\addcontentsline{toc}{section}{Índice de algoritmos}
\renewcommand\listalgorithmname{Índice de algoritmos}
\listofalgorithms

%-----------------------------------------------------------------------------------------------------------------
\end{comment}

% REFERENCIAS


\clearpage
%Let's cite! The Einstein's journal paper \cite{dirac} and the Dirac's 
%book \cite{einstein} are physics related items. 

%\Urlmuskip=0mu plus 1mu\relax
\addcontentsline{toc}{section}{Referencias} 
\printbibliography
 
\end{document}
